\chapter{Foundations and isometric group actions}

This chapter fixes the conventions of Sections 2.1--2.6 of this paper
\cite{bost2026}. All vector spaces are finite dimensional. Subspaces always
carry restricted norms, and quotients always carry quotient norms.

\begin{definition}[The adelic context]
\label{def:adelic-context}
\source{bost-adelic-2026:page-0001,bost-adelic-2026:page-0003,bost-adelic-2026:page-0004}
\lean{ArakelovGeometry.AdelicStructure, ArakelovGeometry.AdelicStructure.IsProper}
\leanok
An adelic curve consists of a field $K$, a measure space
$(\Omega,\mathcal A,\nu)$, and an absolute value $|\cdot|_\omega$ on $K$
for each $\omega\in\Omega$, such that
$\omega\mapsto\log|a|_\omega$ is measurable and integrable for every
$a\in K^\times$. It is \emph{proper} if
\[
  \int_\Omega\log|a|_\omega\,d\nu(\omega)=0
  \qquad(a\in K^\times).
\]
Throughout, $S$ denotes such a proper curve, $K$ is countable of characteristic
zero, and $K_\omega$ denotes the completion of $K$ at $\omega$.
Write $\Omega_\infty$ for the Archimedean places. At each such place we
require an isometric identification of $K_\omega$ with $\mathbb R$ or
$\mathbb C$ equipped with its usual absolute value. Set
$m_\infty=\nu(\Omega_\infty)$.
\end{definition}

\begin{lemma}[A common countable dense subfield]
\label{lem:countable-common-density}
\uses{def:adelic-context}
\source{bost-adelic-2026:page-0004}
\lean{ArakelovGeometry.AdelicStructure.hasCommonCountableDenseSubfield_of_countable}
\leanok
There is a single countable subfield of $K$ whose image is dense in every
$K_\omega$.
\end{lemma}
\begin{proof}
Take the subfield to be $K$ itself. Its image is dense in its completion by
the definition of completion. This works for each place, including a
trivially valued place, and countability is one of our hypotheses on $K$.
\end{proof}

\begin{definition}[Hermitian adelic vector bundles]
\label{def:hermitian-adelic-bundle}
\uses{def:adelic-context}
\source{bost-adelic-2026:page-0004}
\lean{ArakelovGeometry.AdelicVectorBundle, ArakelovGeometry.AdelicVectorBundle.IsHermitian}
\leanok
For a $K$-vector space $E$, put $E_\omega=K_\omega\otimes_K E$ and
$E^\vee=\operatorname{Hom}_K(E,K)$. A norm family $\xi$ assigns a norm
$\|\cdot\|_\omega$ to $E_\omega$. Its dual family is the operator norm
\[
 \|\ell\|_{\omega,*}
   =\sup_{x\in E_\omega\setminus\{0\}}
       \frac{|\ell(x)|_\omega}{\|x\|_\omega}.
\]
The family is \emph{Hermitian} if its norms come from positive definite real
or Hermitian inner products at Archimedean places and are ultrametric at
the remaining places. It is \emph{measurable} if
$\omega\mapsto\|1\otimes s\|_\omega$ is measurable for every $s\in E$.
It is \emph{upper dominated} if, for each nonzero $s\in E$, the function
$\log\|1\otimes s\|_\omega$ is almost everywhere bounded above by an
integrable real function. It is \emph{dominated} if it and its dual are
upper dominated. A Hermitian adelic vector bundle
$\overline E=(E,\xi)$ is a Hermitian dominated norm family for which both
$\xi$ and $\xi^\vee$ are measurable. Its rank is $\dim_K E$.
For the zero vector space the dual norm is zero; equivalently the displayed
supremum is taken in the nonnegative reals with $\sup\varnothing=0$.
\end{definition}

\begin{definition}[Subquotient, tensor, and determinant norms]
\label{def:adelic-induced-norms}
\uses{def:hermitian-adelic-bundle}
\source{bost-adelic-2026:page-0005}
For $A\subset E$, restriction gives the norms on $A_\omega$, and the
norm on $(E/A)_\omega$ is
\[
 \|q\|_\omega=\inf_{x\bmod A_\omega=q}\|x\|_\omega.
\]
For a finite nonempty family $(\overline E_i)_{i\in I}$, its
\emph{orthogonal tensor product} has underlying space
$T=\bigotimes_{i\in I}E_i$. At an Archimedean place its inner product is
specified by
\[
 \left\langle\bigotimes_i x_i,\bigotimes_i y_i\right\rangle_\omega
       =\prod_i\langle x_i,y_i\rangle_\omega.
\]
At a non-Archimedean place its norm is the injective tensor norm
\[
 \|z\|_{\omega,\varepsilon}
   =\sup_{\ell_i\in E_{i,\omega}^\vee\setminus\{0\}}
        \frac{|(\bigotimes_i\ell_i)(z)|_\omega}
             {\prod_i\|\ell_i\|_{\omega,*}}.
\]
The binary formula is that of Section 2.3; the displayed finite formula
uses the same convention for a finite family.
For $r=\dim_K E>0$, the determinant line is $\det E=\bigwedge^r E$.
Its norm at an Archimedean place satisfies
\[
 \|x_1\wedge\cdots\wedge x_r\|_\omega^2
    =\det(\langle x_i,x_j\rangle_\omega)_{i,j=1}^r;
\]
at a non-Archimedean place it is
\[
 \|z\|_\omega
    =\inf_{x_1\wedge\cdots\wedge x_r=z}\prod_{i=1}^r\|x_i\|_\omega.
\]
\end{definition}

The foundational subquotient and norm calculus used below is Proposition 2.3
of this paper: these constructions give Hermitian adelic bundles, and
$(E/A)^\vee\simeq A^\perp\subset E^\vee$ and
$A^\vee\simeq E^\vee/A^\perp$ are isometries. Here
$A^\perp=\{\ell\in E^\vee:\ell(A)=0\}$.
Its measurability hypothesis is supplied by
Lemma~\ref{lem:countable-common-density}. In particular, no discrete set of
places or number-field assumption is imposed.

\begin{definition}[Degree, slope, and semistability]
\label{def:adelic-slope}
\uses{def:adelic-context,def:hermitian-adelic-bundle,def:adelic-induced-norms,lem:countable-common-density}
\source{bost-adelic-2026:page-0005,bost-adelic-2026:page-0006}
\lean{ArakelovGeometry.AdelicVectorBundle.degree, ArakelovGeometry.AdelicVectorBundle.slope, ArakelovGeometry.AdelicVectorBundle.maximalSlope, ArakelovGeometry.AdelicVectorBundle.IsMeanSlopeSemistable}
\leanok
For a Hermitian adelic line $\overline L$ and $s\in L\setminus\{0\}$,
define
\[
 \widehat{\deg}(\overline L)
       =-\int_\Omega\log\|1\otimes s\|_\omega\,d\nu(\omega).
\]
For $\overline E$ of rank $r>0$, define
\[
 \widehat{\deg}(\overline E)=\widehat{\deg}(\det\overline E),\qquad
 \widehat\mu(\overline E)=\frac{\widehat{\deg}(\overline E)}r,\qquad
 \widehat\mu_{\max}(\overline E)
       =\sup_{0\ne A\subset E}\widehat\mu(\overline A).
\]
Set $\widehat{\deg}(0)=0$. A positive-rank bundle is \emph{semistable} if
$\widehat\mu(\overline A)\le\widehat\mu(\overline E)$ for every
nonzero $K$-subspace $A\subset E$. Since $A=E$ is allowed, this is
equivalent to $\widehat\mu_{\max}(\overline E)=\widehat\mu(\overline E)$.
\end{definition}

The line integral is well defined. Indeed, choose $s^\vee$ with
$s^\vee(s)=1$. Upper domination of $s$ and $s^\vee$ bounds
$\log\|s\|_\omega$ above and below by integrable functions, since
$\|s\|_\omega\|s^\vee\|_{\omega,*}=1$. Replacing $s$ by $as$
changes the integral by $-\int\log|a|_\omega\,d\nu=0$.

\begin{lemma}[Exact-sequence degree additivity]
\label{lem:degree-exact-sequence}
\uses{def:adelic-slope,def:adelic-induced-norms}
\source{bost-adelic-2026:page-0006}
\lean{ArakelovGeometry.AdelicVectorBundle.degree_restrict_add_degree_quotient_eq_of_isHermitian}
\leanok
For every subspace $A$ of a Hermitian adelic bundle $\overline E$,
\[
 \widehat{\deg}(\overline E)
   =\widehat{\deg}(\overline A)+\widehat{\deg}(\overline E/\overline A).
\]
This includes $A=0$ and $A=E$.
\end{lemma}
\begin{proof}
Apply the Hermitian degree-additivity theorem, Chen--Moriwaki,
Proposition 4.3.13, equation (4.29), in the form recalled as equation (10)
of this paper. Its subspace norms are restrictions and its quotient
norms are the infima over lifts, exactly as in
Definition~\ref{def:adelic-induced-norms}. The norm families and their
duals are measurable and dominated, and the curve is proper, so its
hypotheses hold. For a zero end the formula reduces to the convention
$\widehat{\deg}(0)=0$. This is an imported foundational degree identity.
\end{proof}

The remaining degree identities used here are the Hermitian dual and tensor
formulas, Chen--Moriwaki, Propositions 4.3.10 and 4.3.19, recalled in
equations (7)--(9) of this paper:
\[
 \widehat{\deg}(\overline E^\vee)=-\widehat{\deg}(\overline E),\qquad
 \widehat\mu(\overline E\otimes\overline F)
     =\widehat\mu(\overline E)+\widehat\mu(\overline F).
\]
These concern mean slopes and do not assert tensor semistability.

\begin{lemma}[Slope duality]
\label{lem:slope-duality}
\uses{def:adelic-slope,def:adelic-induced-norms}
\source{bost-adelic-2026:page-0006}
\lean{ArakelovGeometry.AdelicVectorBundle.slope_add_slope_dual_eq_zero_of_isHermitian, ArakelovGeometry.AdelicVectorBundle.minimalSlope_dual_add_maximalSlope_eq_zero_of_isHermitian, ArakelovGeometry.AdelicVectorBundle.minimalSlope_eq_iInf_slope_quotient}
\leanok
For a positive-rank Hermitian adelic bundle define
\[
 \widehat\mu_{\min}(\overline E)
   =\inf_{A\subsetneq E}\widehat\mu(\overline E/\overline A).
\]
Then
\[
 \widehat\mu(\overline E^\vee)=-\widehat\mu(\overline E),\qquad
 \widehat\mu_{\min}(\overline E^\vee)
       =-\widehat\mu_{\max}(\overline E).
\]
Equivalently, $\widehat\mu_{\min}(\overline E)
=-\widehat\mu_{\max}(\overline E^\vee)$.
\end{lemma}
\begin{proof}
Duality preserves rank, so the dual degree formula proves the mean-slope
identity. Annihilators biject nonzero subspaces $A\subset E$ with
nonzero quotients $E^\vee/A^\perp$. The latter are isometric to
$A^\vee$, and hence have slope $-\widehat\mu(\overline A)$.
Taking the infimum of these negatives gives the negative of the
supremum defining the maximal slope. Applying this to $E^\vee$ and
using the bidual isometry gives the equivalent identity. These are the
Hermitian quotient-mean and duality conventions of equation (12) of
this paper; in particular every quotient carries its actual quotient norm.
\end{proof}

\begin{lemma}[Semistability and duality]
\label{lem:adelic-dual-semistability}
\uses{def:adelic-slope,def:adelic-induced-norms,lem:degree-exact-sequence,lem:slope-duality}
\source{bost-adelic-2026:page-0006}
\lean{ArakelovGeometry.AdelicVectorBundle.isMeanSlopeSemistable_dual_iff_of_isHermitian}
\leanok
A positive-rank Hermitian adelic bundle is semistable if and only if its
dual is semistable.
\end{lemma}
\begin{proof}
Write $r=\dim E$, $a=\dim A$, and $q=r-a$. Degree additivity gives
\[
 q\bigl(\widehat\mu(\overline E/\overline A)
                -\widehat\mu(\overline E)\bigr)
   =a\widehat\mu(\overline E)-\widehat{\deg}(\overline A).
\]
For $a,q>0$ the right side is nonnegative precisely when
$\widehat\mu(\overline A)\le\widehat\mu(\overline E)$. For $A=0$
the quotient inequality is equality; for $A=E$ the subspace inequality is
equality. Thus semistability is equivalent to all nonzero quotient slopes
being at least the slope of $E$.
Every subspace $B\subset E^\vee$ is isometric to the dual of
$E/B^\perp$. The dual degree identity turns that quotient inequality into
$\widehat\mu(\overline B)\le-\widehat\mu(\overline E)
=\widehat\mu(\overline E^\vee)$.
This proves semistability of the dual. Applying the same argument to the
dual, and using the canonical isometry $E\simeq E^{\vee\vee}$, proves
the converse.
\end{proof}

\begin{theorem}[The Hermitian Harder--Narasimhan filtration]
\label{thm:hermitian-hn-filtration}
\uses{def:adelic-slope,lem:degree-exact-sequence,lem:slope-duality}
\source{bost-adelic-2026:page-0007}
\lean{ArakelovGeometry.AdelicVectorBundle.hermitianHNFlag, ArakelovGeometry.AdelicVectorBundle.hermitianHNFlag_strictMono, ArakelovGeometry.AdelicVectorBundle.isMeanSlopeSemistable_hermitianHNFactor, ArakelovGeometry.AdelicVectorBundle.hermitianHNFactor_slope_strictAnti, ArakelovGeometry.AdelicVectorBundle.hermitianHNOrderedSlope_zero_eq_maximalSlope, ArakelovGeometry.AdelicVectorBundle.slope_hermitianHNFactor_last_eq_minimalSlope, ArakelovGeometry.AdelicVectorBundle.sum_hermitianHNFactor_rank_mul_slope_eq_degree, ArakelovGeometry.AdelicVectorBundle.hermitianHNFiltration}
\leanok
Every positive-rank Hermitian adelic bundle $\overline E$ has a flag
\[
 0=E_0\subsetneq E_1\subsetneq\cdots\subsetneq E_\ell=E
 \qquad(\ell\ge1)
\]
whose graded bundles $\overline G_i=\overline E_i/\overline E_{i-1}$
are semistable, with strictly decreasing slopes
$\mu_1>\cdots>\mu_\ell$. Its extreme slopes and degree satisfy
\[
 \mu_1=\widehat\mu_{\max}(\overline E),\qquad
 \mu_\ell=\widehat\mu_{\min}(\overline E),\qquad
 \widehat{\deg}(\overline E)=\sum_{i=1}^\ell\dim_K(G_i)\mu_i.
\]
The associated decreasing real filtration is
\[
 \mathcal H_E^tE=\sum_{\mu_i\ge t}E_i\qquad(t\in\mathbb R).
\]
It is left-continuous, equals $E$ for $t\le\mu_\ell$, and equals $0$
for $t>\mu_1$.
\end{theorem}
\begin{proof}
We invoke the foundational Hermitian HN theorem of Chen--Moriwaki,
Section 4.3.10 and Proposition 4.3.37, as recalled in Theorem 2.5 of
this paper. The properness, Hermitianity, domination, and measurability
assumptions are exactly those in Definitions~\ref{def:adelic-context}
and~\ref{def:hermitian-adelic-bundle}. The construction successively
takes the greatest maximal-slope subspace of the remaining quotient.
That theorem supplies semistability of the resulting nonzero grades
and strict decrease of their slopes. Its first term attains the maximal
slope, and Proposition 4.3.41, as cited in Theorem 2.5, identifies the
last grade's slope with the minimal slope. This is an explicit use of
the imported HN theorem, without any assertion about tensor products.
Repeated use of Lemma~\ref{lem:degree-exact-sequence} gives the displayed
degree sum. Finally, the definition of $\mathcal H_E^tE$ is a finite sum
of nested subspaces selected by $\mu_i\ge t$; it gives the endpoint
values and left-continuity directly.
\end{proof}

\begin{lemma}[Uniqueness of the Hermitian HN filtration]
\label{lem:hn-uniqueness}
\uses{thm:hermitian-hn-filtration}
\source{bost-adelic-2026:page-0007}
\lean{ArakelovGeometry.AdelicVectorBundle.semistableFlag_eq_hermitianHNTerm, ArakelovGeometry.AdelicVectorBundle.hermitianHNFiltration_unique}
\leanok
Any strict flag from $0$ to $E$ with semistable induced graded bundles
and strictly decreasing graded slopes defines the same real filtration
$\mathcal H_E$. Consequently its length, its subspaces, and its graded
slopes agree with the HN flag.
More generally, let $F^tE$ be a decreasing real filtration, equal to $E$
for sufficiently small $t$ and to $0$ for sufficiently large $t$, and
locally constant on an interval immediately to the left of every $t$,
including its right endpoint. Put
$F^{t+}E=\sum_{s>t}F^sE$. If every nonzero
$F^tE/F^{t+}E$ is semistable of slope $t$, then $F^tE=\mathcal H_E^tE$
for all $t$.
\end{lemma}
\begin{proof}
Apply the uniqueness theorem for Hermitian HN flags,
Chen--Moriwaki, Proposition 4.3.38, invoked in Theorem 2.5 of this paper.
The required graded norms are the restriction followed by quotient
norms, and both strictness assumptions are precisely its hypotheses.
It identifies the flags, hence their graded bundles and slopes. The
formula defining the real filtration then identifies the filtrations
at every $t$. Conversely its distinct nonzero values recover the flag
and its jump parameters recover the graded slopes.
For the real-filtration assertion, its distinct values form a finite
chain: every strict inclusion increases dimension. The endpoint and
left-continuity hypotheses therefore express it as a finite flag
with decreasing real jump parameters, including each jump at its
parameter. Its nonzero graded quotients are the successive quotients
of this flag. The assumed equality of their slopes with the jump
parameters gives strict decrease. Apply the flag uniqueness just proved
to identify the flag and all its jump parameters, and thus the original
real filtration.
\end{proof}

\begin{theorem}[The maximal-slope part of the Hermitian HN theorem]
\label{thm:largest-destabilizing-subspace}
\uses{thm:hermitian-hn-filtration}
\source{bost-adelic-2026:page-0006,bost-adelic-2026:page-0007}
\lean{ArakelovGeometry.AdelicVectorBundle.exists_unique_largest_maximalSlope_subspace}
\leanok
For every positive-rank Hermitian adelic bundle $\overline E$, there is a
unique nonzero subspace $D_E\subset E$ such that
\[
 \widehat\mu(\overline D_E)=\widehat\mu_{\max}(\overline E),\qquad
 \widehat\mu(\overline A)=\widehat\mu_{\max}(\overline E)
       \Longrightarrow A\subset D_E\quad(0\ne A\subset E).
\]
In particular the maximal slope is a finite real number.
\end{theorem}
\begin{proof}
We use the Hermitian Harder--Narasimhan theorem recalled in Theorem 2.5
of this paper, including its description of the first nonzero term as
the greatest subspace of maximal slope. Apply that theorem to
$\overline E$ and take $D_E$ to be this first term. It is nonzero,
has slope equal to the maximal slope, and contains every subspace of
that slope. These are exactly the asserted properties. If $D'_E$ also
has them, each of $D_E,D'_E$ is contained in the other, proving uniqueness.
The slope of the positive-rank bundle $\overline D_E$ is real by the
integrability in the definition of degree, proving finiteness.
\end{proof}

This is the foundational HN input used here; only its first term is needed.
The subsequent group argument does not require a tensor-product HN theorem.

\begin{lemma}[Isometries preserve the largest destabilizing subspace]
\label{lem:isometry-preserves-destabilizer}
\uses{thm:largest-destabilizing-subspace,def:adelic-induced-norms}
\source{bost-adelic-2026:page-0007}
\lean{ArakelovGeometry.AdelicVectorBundle.map_destabilizingSubspace_eq_of_completed_isometry}
\leanok
Let $u:E\xrightarrow{\sim}F$ be a $K$-linear isomorphism between
positive-rank Hermitian adelic bundles such that every
$u_\omega:E_\omega\to F_\omega$ is an isometry. Then $u(D_E)=D_F$.
\end{lemma}
\begin{proof}
For each subspace $A\subset E$, the restriction
$A_\omega\to u(A)_\omega$ is an isometry. The determinant norm formulas
show that its determinant is an isometry as well: at Archimedean places
the Gram matrices agree, and at other places $u_\omega$ bijects the
wedge decompositions in the defining infimum without changing their
products of norms. Thus $A$ and $u(A)$ have equal degrees and ranks,
and hence equal slopes when nonzero. Applying this to $u$ and $u^{-1}$
identifies the two maximal slopes. Consequently $u(D_E)$ has maximal
slope and lies in $D_F$. The same argument applied to $u^{-1}$ yields
$u^{-1}(D_F)\subset D_E$, proving equality.
\end{proof}

\begin{lemma}[Irreducible isometric group actions]
\label{lem:group-semistability}
\uses{lem:isometry-preserves-destabilizer,def:adelic-slope}
\source{bost-adelic-2026:page-0008}
\lean{ArakelovGeometry.AdelicVectorBundle.isMeanSlopeSemistable_of_isIrreducible}
\leanok
Let a group $G$ act $K$-linearly on a positive-rank Hermitian adelic
bundle $\overline E$. Suppose its representation is irreducible over $K$
and every completed action is isometric:
\[
 \|g x\|_\omega=\|x\|_\omega
 \qquad(g\in G,\ \omega\in\Omega,\ x\in E_\omega).
\]
Then $\overline E$ is semistable. In particular this holds for a finite group.
\end{lemma}
\begin{proof}
For every $g\in G$, Lemma~\ref{lem:isometry-preserves-destabilizer}
gives $gD_E=D_E$. Hence $D_E$ is a nonzero invariant subspace.
Irreducibility forces $D_E=E$. The defining property of $D_E$ then gives
$\widehat\mu_{\max}(\overline E)=\widehat\mu(\overline E)$.
\end{proof}

The proof of the first assertion of Lemma 2.6 of this paper uses no
finiteness of $G$; the statement above records this harmless generalization.

\begin{lemma}[The action spans all endomorphisms]
\label{lem:irreducible-action-span}
\lean{BostConjecture.Arithmetic.span_representation_eq_top}
\leanok
Let $L$ be an algebraically closed field and $G$ a group. Let $V$ be a nonzero
finite dimensional vector space over $L$, with an irreducible action $\rho$
of $G$. Then
\[
 \operatorname{span}_L\{\rho(g):g\in G\}
       =\operatorname{End}_L(V).
\]
\end{lemma}
\begin{proof}
The action extends to the group algebra $L[G]$, and its image is exactly
the displayed span. Schur's lemma says that every nonzero
$L[G]$-endomorphism of $V$ is invertible. If $a$ is such an endomorphism,
algebraic closedness gives an eigenvalue $\lambda\in L$; the endomorphism
$a-\lambda\operatorname{id}$ has a nonzero kernel and therefore must be
zero. Thus $\operatorname{End}_{L[G]}(V)=L\operatorname{id}$.
The finite-dimensional Jacobson density theorem says that the image of
$L[G]$ is the endomorphism algebra over this commutant. It is consequently
$\operatorname{End}_L(V)$, as required.
\end{proof}

\begin{lemma}[Irreducibility of the independent product action]
\label{lem:independent-product-irreducible}
\uses{lem:irreducible-action-span}
\source{bost-adelic-2026:page-0008}
\lean{BostConjecture.Arithmetic.isIrreducible_independentPi_of_representationBaseChange}
\leanok
Let $I$ be a finite nonempty set, let $G_i$ be groups, and let $E_i$ be
nonzero finite-dimensional $K$-representations. Suppose there is an
algebraically closed extension $L/K$ such that each $L\otimes_K E_i$ is
irreducible. Then $T=\bigotimes_{i\in I}E_i$ is irreducible for the
\emph{independent product action}
\[
 (g_i)_{i\in I}\left(\bigotimes_i x_i\right)=\bigotimes_i(g_i x_i)
 \qquad\left((g_i)\in\prod_{i\in I}G_i\right).
\]
In particular the hypothesis holds when every $E_i$ is absolutely
irreducible, meaning that its scalar extension to an algebraic closure
of $K$ is irreducible.
\end{lemma}
\begin{proof}
Put $V_i=L\otimes_K E_i$. The preceding lemma says that each factor
action spans $\operatorname{End}_L(V_i)$. The natural map
\[
 \bigotimes_i\operatorname{End}_L(V_i)
       \longrightarrow\operatorname{End}_L\left(\bigotimes_i V_i\right)
\]
is an isomorphism: after choosing bases, tensor products of matrix units
map bijectively to the matrix units in the tensor basis. Expanding each
factor in action operators proves that the independent product action
spans the endomorphisms of $\bigotimes_i V_i$. A nonzero invariant
subspace is therefore stable under every endomorphism. Given its nonzero
vector $v$ and any vector $w$, choose a linear functional $f$ with
$f(v)=1$; the endomorphism $x\mapsto f(x)w$ sends $v$ to $w$.
The subspace is thus the whole tensor product, which is nonzero by its
tensor basis.

The natural isomorphism $L\otimes_K T\simeq\bigotimes_i V_i$ intertwines
the actions, as is seen on pure tensors. If $0\ne A\subsetneq T$ were
invariant, then $L\otimes_K A$ would be a nonzero proper invariant
subspace: scalar extension preserves dimensions, and
$0<\dim_K A<\dim_K T$. This contradicts the irreducibility just proved.
\end{proof}

\begin{lemma}[Independent actions preserve the tensor norms]
\label{lem:independent-product-isometry}
\uses{def:adelic-induced-norms}
\source{bost-adelic-2026:page-0008}
\lean{ArakelovGeometry.NormFamily.orthogonalTensor_norm_independentPiRepresentation}
\leanok
Let $I$ be a finite nonempty set and, for each $i\in I$, let $G_i$ act
$K$-linearly on a positive-rank Hermitian adelic bundle $\overline E_i$.
Suppose each $G_i$ acts by isometries on every completed fibre of
$\overline E_i$. Then the independent product action preserves every
completed orthogonal tensor norm of $\bigotimes_i\overline E_i$.
\end{lemma}
\begin{proof}
At an Archimedean place, a linear isometry preserves the inner product
by polarization. The tensor inner product formula therefore gives
equality on pairs of pure tensors; sesquilinearity gives equality on
arbitrary pairs, and hence equality of norms.
At a non-Archimedean place, an invertible isometry $u_i$ induces a
bijection $\ell_i\mapsto\ell_i\circ u_i$ on nonzero linear functionals,
preserving dual norms, since the change of variables $y=u_i x$ bijects
the vectors in the defining supremum. Substituting these functionals
in the formula for the injective tensor norm proves
$\|(\bigotimes_i u_i)z\|_\varepsilon=\|z\|_\varepsilon$.
Apply this with $u_i=g_i$ on each completed fibre.
\end{proof}

\begin{lemma}[Semistability for independent group actions]
\label{lem:independent-group-semistability}
\uses{lem:group-semistability,lem:independent-product-irreducible,lem:independent-product-isometry,lem:adelic-dual-semistability}
\source{bost-adelic-2026:page-0008}
\lean{ArakelovGeometry.AdelicVectorBundle.isMeanSlopeSemistable_orthogonalTensor_and_dual_of_isIrreducible_baseChange}
\leanok
Let $(G_i,\overline E_i)_{i\in I}$ be a finite nonempty family of
positive-rank Hermitian adelic bundles with group actions preserving
every completed norm. Suppose that, over one common algebraically
closed extension $L/K$, all factor representations become irreducible.
Then the orthogonal tensor bundle $\bigotimes_i\overline E_i$ and its
dual are semistable. In particular this proves the second assertion of
Lemma 2.6 of this paper for absolutely irreducible representations of
finite groups.
\end{lemma}
\begin{proof}
The tensor bundle has positive rank $\prod_i\dim_K E_i$.
Lemma~\ref{lem:independent-product-irreducible} proves irreducibility
of its independent product representation, and
Lemma~\ref{lem:independent-product-isometry} proves preservation of
its completed norms. Apply Lemma~\ref{lem:group-semistability} to obtain
semistability. The dual is semistable by
Lemma~\ref{lem:adelic-dual-semistability}.
\end{proof}

\begin{corollary}[Line factors in an irreducible isometric action]
\label{cor:group-line-twist-semistability}
\uses{lem:group-semistability,lem:independent-product-irreducible,lem:independent-product-isometry,lem:adelic-dual-semistability}
\lean{ArakelovGeometry.AdelicVectorBundle.isMeanSlopeSemistable_lineTensor_and_dual_of_isIrreducible, BostConjecture.Arithmetic.isIrreducible_baseChange_trivialLineTensor}
Let $G$ act irreducibly on a positive-rank Hermitian adelic bundle
$\overline V$ by isometries of every completed fibre, and let
$\overline L$ be any Hermitian adelic line bundle. Then
$\overline L\otimes\overline V$ and its dual are semistable.
The action on $L\otimes_K V$ is trivial on $L$ and the given action on $V$;
it stays irreducible after every field extension over which the action on
$V$ is irreducible.
In particular this applies with
$\overline V=\bigotimes_i\overline E_i$ under the hypotheses of
Lemma~\ref{lem:independent-group-semistability}.
\end{corollary}
\begin{proof}
Give $L$ the trivial $G$-action. Choose $0\ne l\in L$. The isomorphism
$V\to L\otimes_K V$, $v\mapsto l\otimes v$, intertwines the original
action with the action $g(l'\otimes v)=l'\otimes gv$.
This action is therefore irreducible. It preserves the completed tensor
norms: the inner product argument and the change of dual functionals
in Lemma~\ref{lem:independent-product-isometry} apply to
$\operatorname{id}_{L_\omega}\otimes g_\omega$.
Lemma~\ref{lem:group-semistability} proves semistability, and
Lemma~\ref{lem:adelic-dual-semistability} proves it for the dual.
Scalar-extending the displayed intertwining isomorphism proves the assertion
about any field extension.
For the indicated finite tensor bundle, irreducibility and completed
isometry follow respectively from
Lemmas~\ref{lem:independent-product-irreducible}
and~\ref{lem:independent-product-isometry}.
\end{proof}
